%%
% 第一章：项目背景与实验目标
%%

\chapter{项目背景与实验目标}
\label{chap:background}

\section{PageIndex项目概述}

PageIndex是一个基于LangChain框架的智能文档问答系统，其核心功能是对索引文档进行语义搜索并生成基于文档内容的回答。该系统采用PostgreSQL作为后端存储，支持对大量文档进行高效管理和检索。从技术架构上看， PageIndex系统主要包含三个核心模块：文档索引模块负责将PDF、Markdown等格式的文档解析为结构化的树形数据并存储到数据库中； 文档搜索模块接收用户的自然语言查询,首先通过大语言模型从文档目录中选择最相关的文档，然后在这些文档的结构树中搜索相关节点，最后将搜索结果组装成上下文并生成最终答案; LangChain集成模块将整个搜索流程封装为LangChain工具，使其可以方便地集成到AI Agent应用中。

PageIndex项目的代码结构清晰地体现了这一设计理念。核心文件\texttt{search.py}（约400行）实现了完整的搜索流程， 包括数据模型定义、文档选择、并发搜索、上下文构建等关键功能。该文件定义了\texttt{CatalogEntry}、\texttt{RelevantNode}、\texttt{DocumentSearchResult}、\texttt{PageIndexSearchResult}等数据类来封装中间结果， 并通过\texttt{run\_tree\_search}函数协调整个搜索流程。另一个核心文件\texttt{page\_index.py}（约2000行）实现了文档解析和索引的底层逻辑， 包含大量的异步处理、并发控制和条件分支， 这些复杂的控制流结构使其成为代码覆盖率分析的理想测试对象。

从程序结构角度分析, PageIndex项目代码包含了丰富的控制流特征。在顺序结构方面, 代码包含大量的数据处理流程, 如JSON解析、数据转换、结果格式化等。在选择结构方面, 代码中广泛使用了if/elif/else语句进行条件判断, 包括参数验证（如检查\texttt{doc\_top\_k}是否大于0）、类型检查（如判断payload是否为dict或list）、状态判断（如判断搜索是否成功）等。在循环结构方面, 代码使用了for循环遍历文档列表和节点树, 使用while循环进行栈式遍历, 并通过asyncio实现并发控制。此外, 代码还包含异常处理结构, 对数据库操作、API调用等可能失败的操作进行异常捕获。这些丰富的控制流结构为执行路径追踪和覆盖率分析提供了充分的测试场景。

\section{实验目标与要求}

本次课程作业的核心目标是设计并实现一个能够记录Python程序执行路径和覆盖率的工具。根据作业要求, 该工具需要满足以下功能: 能够记录Python程序的执行路径和覆盖率; 能够分析基本程序常见结构流, 包括选择、顺序与循环; 能够记录程序运行时的执行路径; 能够生成描述程序运行执行分支的报告。同时, 代码需要有必要的注释和说明使用文档, 程序结构需要清晰且具有一定的可扩展性。

为了实现这些目标, 我们设计了PathTracer工具, 它是一个专门用于追踪Python程序执行路径和计算代码覆盖率的工具。与现有的coverage.py等工具相比, PathTracer的独特之处在于它同时结合了静态分析和动态追踪两种技术: 通过AST静态分析识别源代码中的所有控制流分支点, 通过运行时追踪记录实际执行的代码路径, 然后通过对比两者来计算覆盖率。此外, PathTracer还具备函数调用树追踪功能, 能够记录函数调用的参数和返回值, 这对于理解程序的执行行为具有重要价值。

\section{技术背景}

在实现PathTracer工具之前, 我们需要了解几个关键技术。首先是Python的\texttt{sys.settrace}机制, 这是Python解释器提供的运行时追踪接口。当通过\texttt{sys.settrace(trace\_callback)}注册一个回调函数后, Python解释器会在程序执行过程中每当发生特定事件时调用该回调。这些事件包括: 'call'事件在函数被调用时触发; 'line'事件在即将执行新的一行代码时触发; 'return'事件在函数返回时触发; 'exception'事件在发生异常时触发。通过处理这些事件, 我们可以精确记录程序的执行轨迹。需要注意的是, \texttt{sys.settrace}会带来显著的性能开销, 程序执行速度可能降低10到100倍, 因此该技术更适合教学演示、调试分析等场景, 而不适合在生产环境中常驻使用。

第二个关键技术是Python的\texttt{ast}模块, 它可以将源代码解析为抽象语法树。AST是源代码的树形表示, 其中每个节点代表代码中的一个结构。通过遍历AST, 我们可以识别源代码中的所有控制流结构, 包括\texttt{ast.If}（if条件语句）、\texttt{ast.For}（for循环）、\texttt{ast.While}（while循环）、\texttt{ast.ExceptHandler}（异常处理器）等。这种静态分析方法使我们能够在不执行程序的情况下获取代码的完整结构信息, 这是计算代码覆盖率的基础——我们需要知道代码中有哪些分支点, 然后才能判断哪些分支在运行时被执行了。

代码覆盖率是软件测试领域的重要概念, 用于衡量测试用例对被测代码的覆盖程度。常见的覆盖率类型包括: 语句覆盖测量被执行的代码语句比例; 分支覆盖测量控制流分支的执行比例, 如if语句的true/false两个分支是否都被执行; 路径覆盖测量程序执行路径的覆盖比例。本实验主要关注分支覆盖, 即识别程序中的控制流分支点, 并统计在运行时被触发的分支数量。覆盖率公式为: 覆盖率 = (已执行分支数 / 总分支数) × 100\%。需要注意的是, 高覆盖率本身不是目的, 重要的是覆盖的有效性——如果未覆盖的分支主要是防御性代码（如空值检查、异常处理）或互斥逻辑（不可能同时执行的分支）, 那么即使覆盖率不是100\%也是可以接受的。
